% =============================================================================
% Capítulo 3 — X binário e Z contínuo: MTE
% Guia CLEAR de Variáveis Instrumentais — FGV EESP CLEAR
% =============================================================================

\chapter{Quando o instrumento é contínuo: quem está na margem?}
\label{cap:mte}

No capítulo anterior, o instrumento era binário: a pessoa ganhou ou não ganhou a loteria. Isso criou dois grupos bem definidos, sorteados e não sorteados, e a razão de Wald estimou o efeito do tratamento para os compliers de forma direta.

Mas muitos instrumentos relevantes na prática não são binários. A distância até uma universidade não é zero ou um: é uma variável contínua. A intensidade de um incentivo fiscal varia de forma contínua. O custo de transporte para chegar a um centro de saúde assume muitos valores diferentes. Nesses casos, o instrumento cria variação \textit{gradual} na probabilidade de tratamento: quanto maior o incentivo (ou menor o custo), maior a probabilidade de participar.

Quando o instrumento é contínuo e o tratamento é binário, a lógica de identificação se enriquece de forma importante. Em vez de um único grupo de compliers, há algo como uma sequência de grupos: à medida que o instrumento aumenta, ele vai "convencendo" novas pessoas a entrar no tratamento, uma margem de cada vez. O Efeito de Tratamento Marginal (MTE) é justamente o nome para o efeito do tratamento sobre as pessoas que estão exatamente nessa margem.

Este capítulo desenvolve essa ideia de forma intuitiva. A matemática fica em segundo plano; o objetivo é entender \textit{o que} o MTE mede e \textit{por que} ele importa para interpretar estimativas de IV.

% =============================================================================
\section{Pessoas diferentes respondem ao instrumento de formas diferentes}
\label{sec:heterogeneidade_selecao}

Vamos mudar o exemplo para tornar a intuição mais clara. Suponha que queremos estimar o efeito de cursar o ensino superior sobre a renda futura. O tratamento $X_i$ é binário: a pessoa cursou ou não a faculdade. O instrumento é a distância da residência até a universidade mais próxima --- quanto maior a distância, maior o custo de cursar a faculdade (transporte, moradia, tempo de deslocamento) e, portanto, menor a probabilidade de cursá-la.

\citet{Card1995} explorou exatamente essa variação. A ideia é que estudantes que por acaso moram perto de uma universidade boa têm uma "vantagem acidental" que aumenta a chance de cursarem o ensino superior, independentemente de suas características pessoais. Essa variação acidental, determinada por onde a família morou, e não por escolha do estudante, é a variação exógena que o instrumento captura.

Mas pense no seguinte: a distância não age de forma idêntica em todo mundo. Para alguém muito motivado, que tinha planos firmes de ir à faculdade de qualquer forma, morar um pouco mais longe não muda nada: ela vai de qualquer jeito. Para alguém que estava na dúvida, que talvez fosse se a logística fosse fácil, mas desistiria diante de qualquer obstáculo a mais: uma distância maior pode ser o fator decisivo. E para alguém que claramente não ia à faculdade independentemente de qualquer obstáculo, a distância também não importa.

Isso significa que à medida que a distância aumenta (ou seja, o instrumento "fica mais adverso"), as pessoas que ``saem'' do tratamento são justamente aquelas que estavam na margem, ou seja, as que precisavam de pouco incentivo extra para desistir. As pessoas com alta determinação continuam cursando a faculdade mesmo com distâncias maiores. Apenas as que estavam em dúvida são afetadas.

O ponto crucial é que esses grupos marginais podem ter efeitos do tratamento sistematicamente diferentes. Pense: quem precisa de muita "pressão" para desistir da faculdade é provavelmente muito motivado e talvez já esperasse colher altos retornos da educação. Quem desiste facilmente talvez tivesse retornos esperados menores. Se isso for verdade, o efeito da faculdade sobre a renda pode ser maior para os mais determinados e menor para os menos determinados.

Essa é a essência da heterogeneidade que o Efeito de Tratamento Marginal captura.

% =============================================================================
\section{A ideia do Efeito de Tratamento Marginal}
\label{sec:mte_definicao}

Para formalizar essa intuição, é útil pensar em cada pessoa como tendo uma certa "resistência ao tratamento": um número que resume todos os fatores não observados que determinam se ela vai ou não tratar quando o instrumento tem um dado valor. Chamamos essa resistência de $U_i$.

Pessoas com $U_i$ pequeno têm baixa resistência: entrariam no tratamento mesmo com pouco incentivo. Pessoas com $U_i$ grande têm alta resistência: precisariam de muito incentivo para entrar, ou não entrariam de jeito nenhum.

A probabilidade de tratamento para um dado valor do instrumento, $p(z) = \mathbb{P}(X=1 \mid Z=z)$, funciona como um "limiar": a pessoa entra no tratamento se sua resistência $U_i$ for menor do que $p(z)$. Quando o instrumento fica mais favorável (maior incentivo, menor distância), $p(z)$ aumenta e o limiar sobe, "puxando" para dentro do tratamento as pessoas com resistências um pouco maiores.

O \textbf{Efeito de Tratamento Marginal} (MTE) é definido como o efeito causal do tratamento para as pessoas que estão exatamente na margem, aquelas com resistência igual ao limiar $u$:
\begin{equation}
\text{MTE}(u) = \E[Y_i(1) - Y_i(0) \mid U_i = u]
\label{eq:mte}
\end{equation}

Ao estimar o MTE para diferentes valores de $u$, traçamos um perfil de como o efeito do tratamento varia ao longo da distribuição de resistência. Se o MTE é decrescente em $u$ (o que acontece quando há seleção positiva para o tratamento), quem entra no tratamento voluntariamente (baixo $U$) tem retornos maiores do que quem só entra quando muito incentivado (alto $U$).

\begin{exemplobox}[MTE e retornos à educação]
\citet{HeckmanVytlacil2005} estimaram o MTE do ensino superior sobre a renda usando distância até a faculdade como instrumento. Eles encontraram que o MTE tende a ser decrescente na resistência ao tratamento: pessoas que teriam ido à faculdade de qualquer jeito (baixo $U$) têm retornos maiores do que pessoas que só foram convencidas pela proximidade geográfica (alto $U$).

Isso faz sentido econômico. Quem tem retornos esperados altos tem mais incentivo para ir à faculdade independentemente dos obstáculos. Quem precisa de condições excepcionalmente favoráveis (morar do lado da universidade) talvez seja menos confiante nos retornos, e de fato acabe com retornos menores.

A implicação prática: um instrumento que afeta principalmente estudantes que moram longe (e que são convencidos a ir à faculdade quando a distância diminui) vai estimar o efeito da faculdade para um grupo com retornos relativamente baixos. Um experimento que convença aleatoriamente pessoas de alta resistência pode estimar um efeito bem diferente do ATE.
\end{exemplobox}

% =============================================================================
\section{Por que diferentes instrumentos dão resultados diferentes}
\label{sec:mte_late}

Uma das implicações mais importantes do framework de MTE é que instrumentos diferentes identificam médias do MTE em intervalos diferentes e, portanto, podem dar estimativas de IV muito distintas mesmo que todos sejam válidos.

Pense em dois instrumentos para o efeito da faculdade sobre a renda:
\begin{enumerate}[leftmargin=*]
  \item Distância até a universidade mais próxima: afeta principalmente pessoas que estavam ``na dúvida'' por razões logísticas. Elas tendem a ter resistência moderada (não muito alta, não muito baixa).
  \item Existência de uma política de cotas que aumentou o acesso para grupos específicos: afeta pessoas que antes não poderiam acessar o ensino superior por razões estruturais. Elas podem ter alta resistência não por falta de motivação, mas por restrições externas.
\end{enumerate}

Cada instrumento move um grupo diferente de pessoas para dentro do tratamento. O LATE identificado por cada instrumento é uma média do MTE sobre o intervalo de resistência que o instrumento cobre. Se o MTE varia ao longo desse intervalo, os dois LATEs serão diferentes, não porque um dos instrumentos está errado, mas porque eles identificam efeitos para grupos diferentes.

Isso tem uma implicação importante para a interpretação de resultados. Quando dois estudos usam IV para estimar o mesmo parâmetro e chegam a números diferentes, a primeira pergunta não deveria ser "qual dos dois está errado?", mas sim "o instrumento de cada estudo afeta qual grupo de pessoas?". A diferença nas estimativas pode ser perfeitamente coerente, refletindo heterogeneidade genuína no efeito do tratamento.

\begin{atencaobox}[Dois instrumentos válidos podem dar resultados diferentes]
Se você usa distância até a universidade como instrumento e encontra um LATE de 15\% de aumento na renda, enquanto outro estudo usa cotas e encontra 8\%, isso não é contradição. Os dois instrumentos podem ser igualmente válidos mas estar identificando efeitos em margens de seleção diferentes. O framework de MTE fornece a linguagem para reconciliar essas diferenças: compare em que trecho da distribuição de resistência cada instrumento está operando.
\end{atencaobox}

% =============================================================================
\section{O MTE como unificador de parâmetros}
\label{sec:mte_unificador}

Uma das contribuições mais elegantes do trabalho de \citet{HeckmanVytlacil2005} é mostrar que os parâmetros que usualmente nos interessam (ATE, ATT e LATE) são todos médias ponderadas do MTE, com pesos diferentes.

O ATE é a média do MTE ao longo de toda a distribuição de resistência, com pesos iguais para todos os valores de $u$. Ele pergunta: qual seria o efeito do tratamento se escolhêssemos aleatoriamente alguém da população?

O ATT é uma média do MTE, mas ponderada pela probabilidade de que uma pessoa com dado $U$ seja tratada. Como pessoas com menor resistência têm mais chance de ser tratadas, o ATT dá mais peso às partes da curva de MTE com $u$ pequeno, onde o efeito tende a ser maior quando há seleção positiva.

O LATE para um instrumento específico é uma média do MTE sobre o intervalo de resistência coberto por aquele instrumento.

Essa estrutura unificada é conceitualmente poderosa porque nos diz que todos esses parâmetros são consistentes entre si: são simplesmente médias do mesmo objeto (o MTE) com pesos diferentes. Quando encontramos que ATT $>$ LATE $>$ ATE (como sugere o padrão típico com seleção positiva), isso reflete os diferentes pesos que cada parâmetro atribui às diferentes partes da curva de MTE.

% =============================================================================
\section{Estimação e limitações práticas}
\label{sec:mte_estimacao}

A estimação do MTE é mais exigente do que a do LATE binário. O processo tem duas etapas: primeiro estima-se como o instrumento afeta a propensão ao tratamento ($p(z)$); depois usa-se a variação nessa propensão para recuperar como o resultado varia ao longo das margens de seleção.

Na prática, isso equivale a estimar como $\E[Y \mid p(Z) = p]$ varia com $p$, e a derivada dessa função em relação a $p$ é o MTE avaliado em $u = p$. Para que essa estimação funcione, precisa haver variação suficiente no instrumento em torno de cada ponto de $u$ que se quer estimar: caso contrário, o MTE não é identificado naquele ponto.

A estimação da função completa do MTE geralmente requer métodos não paramétricos ou semiparamétricos, e a inferência precisa levar em conta que a propensão estimada no primeiro estágio carrega incerteza que se propaga para o segundo. Isso torna o MTE um objeto mais exigente computacionalmente do que o simples LATE.

Para a maioria das aplicações aplicadas, o pesquisador não precisa estimar toda a curva de MTE. É suficiente: (i) estar ciente de que o LATE identificado pelo instrumento disponível pode diferir do ATE; (ii) pensar em qual trecho da distribuição de resistência o instrumento opera; e (iii) discutir explicitamente quem são os compliers afetados pelo instrumento e se o efeito para eles é provavelmente maior, menor ou parecido com o efeito médio para a população.

\begin{notabox}[Pontos-chave do capítulo]
\begin{itemize}[leftmargin=*]
  \item Quando o instrumento é contínuo, ele "convence" pessoas diferentes em margens diferentes. O MTE captura o efeito do tratamento para as pessoas na margem de indiferença para cada valor do instrumento.
  \item Diferentes instrumentos identificam médias do MTE em intervalos diferentes de resistência. Estimativas de IV diferentes para o mesmo tratamento podem refletir heterogeneidade real, não contradição.
  \item ATE, ATT e LATE são todos médias ponderadas do MTE com pesos diferentes. Quando o MTE é decrescente na resistência (seleção positiva), espera-se ATT $>$ LATE $>$ ATE.
  \item O MTE exige mais para ser estimado do que o LATE binário, mas mesmo sem estimá-lo formalmente, o conceito é útil para interpretar e comparar estimativas de IV produzidas por instrumentos diferentes.
\end{itemize}
\end{notabox}
